\section*{Введение}
\addcontentsline{toc}{section}{Введение}

Способы снижения энергии нейросетевых вычислений развивались вместе
с моделями нейронных сетей, численными методами и архитектурой
вычислительных устройств. Применение нейронных сетей в задачах классификации,
регрессии, восстановления и генерации данных предъявляет требования
к производительности, объёму памяти и энергетическим затратам.
Характеристики вычислительной системы зависят от состава операций, организации
обмена данными и используемой арифметики. В обзоре В.~Сзе
и соавторов 2017 года обобщены способы сокращения затрат
на арифметические операции и обмен данными, а также подходы
к аппаратной реализации нейронных сетей~\cite{sze2017}.

Инференсом называется применение обученной нейронной сети к входным данным
для получения ответа модели. При обучении параметры настраиваются
по обучающим данным; в градиентных методах для этого вычисляются
градиенты функции потерь. Различия между этими процессами
учитываются при разработке специализированной аппаратуры: устройство инференса
может быть рассчитано на многократное выполнение вычислений
с заранее подготовленными параметрами при заданных ограничениях
по задержке, пропускной способности
и качеству результата.

Основные способы снижения энергетических затрат включают сокращение
числа операций, уменьшение объёма хранимых и передаваемых данных,
использование численных форматов пониженной разрядности и специализацию
вычислительных средств. На системном уровне результат зависит
от организации параллелизма, размещения данных и управления исполнением.
Перечисленные способы затрагивают различные составляющие затрат
и требуют совместной оценки с качеством решения целевой задачи.

Эти методы опираются на математические модели нейронных сетей, теорию
параллельных вычислений, архитектуру памяти и численный анализ. Каждая из этих
областей имела свои задачи и методы: исследователи изучали
обучение, распределение операций, передачу данных и погрешности вычислений.
Рост сложности нейросетевых моделей и ограничения аппаратуры потребовали
совместного применения этих методов. Разработка специализированных ускорителей
объединила исследования алгоритмов с проектированием вычислительных устройств.

Ранние способы аппаратного выполнения математических операций
разрабатывались в аналоговых машинах 1930-х годов.
В 1943 году У.~Маккаллох и У.~Питтс предложили логическую модель
нейрона. В последующих исследованиях развивались методы обучения,
средства параллельной обработки и памяти, способы сокращения
вычислительной сложности сетей и специализированная аппаратура.
В работах 2010-х и начала 2020-х годов выбор модели и численного
представления рассматривался совместно с организацией вычислительной
системы для сокращения затрат при выполнении модели.
Эти процессы частично совпадали во времени; новые методы дополняли ранее
разработанные способы вычислений.

Изменение требований к вычислениям рассматривается на примерах
свёрточных сетей, визуальных трансформеров и диффузионных моделей.
Для этих моделей анализа и синтеза изображений различаются состав операций,
порядок их выполнения и требования к вычислительной системе.

В работе использованы научные публикации с описанием формальных моделей,
алгоритмов обучения, аппаратных систем и методов их оценки. При сравнении
публикаций различаются теоретические результаты, оценки проектируемых
устройств и измерения на изготовленной аппаратуре. Для численных методов учитываются
формат операндов, правило округления, организация масштабирования
и критерий допустимого изменения результата. При сравнении архитектур
учитываются части системы, для которых получены показатели,
состав вычислительной нагрузки и условия исполнения.

Вместе с методами вычислений развивались критерии их оценки. Показатели числа операций и пиковой
производительности дополнялись характеристиками обмена данными,
энергии полного исполнения и качества решения целевой задачи.
Сравнение этих критериев показывает, какие задачи решали алгоритмические,
численные и аппаратные методы на разных этапах. Сходство общих принципов при
этом не означает одинакового устройства систем или одинаковых условий
получения результата.
